\documentclass[11pt]{article}
\usepackage[margin=1.1in]{geometry}
\usepackage{amsmath,amssymb,amsthm}
\usepackage{booktabs}
\usepackage{longtable}
\usepackage{url}
\usepackage[hidelinks]{hyperref}

\newtheorem{theorem}{Theorem}
\newtheorem{lemma}{Lemma}
\newtheorem{conjecture}{Conjecture}
\newtheorem{claim}{Claim}

\newcommand{\Nout}{N^{+}}
\newcommand{\Nsec}{N^{++}}

\title{Exhaustive verification of Seymour's Second Neighborhood Conjecture
for oriented graphs on at most 14 vertices\thanks{Working note, not
peer reviewed. The computations and this text were produced by an automated
agent pipeline (Claude Fable 5, Atris workspace of Keshav Rao) on a single
shared Apple Silicon laptop, 2026-07-22/23. Section~\ref{sec:audit} lists
exactly what a human reviewer must audit before any of this is submitted
anywhere. No claim in this note should be trusted beyond the evidence
described in it.}}
\author{Automated verification pipeline (Atris workspace)}
\date{July 23, 2026}

\begin{document}
\maketitle

\begin{abstract}
Seymour's Second Neighborhood Conjecture states that every oriented graph
(finite simple digraph with no loops and no 2-cycles) contains a vertex $v$
with $|\Nsec(v)| \ge |\Nout(v)|$, where $\Nsec(v)$ is the set of vertices at
directed distance exactly two from $v$. We report an exhaustive computational
verification: no counterexample exists on $n \le 14$ vertices. The result
rests on 33 unsatisfiable SAT instances solved by CaDiCaL 3.0.1 with DRAT
proof logging, every proof checked by drat-trim; the instances cover the
search space through a minimum out-degree sharding scheme whose covering and
symmetry-breaking arguments are short pen-and-paper proofs reproduced here.
(The full artifact set spans $n \le 18$ and 63 CNF instances, of which 59
refutations are complete; the 33 form the $n \le 14$ claim.)
Independently of the SAT chain, all oriented graphs on $n \le 8$ vertices
(575{,}016{,}219 at $n=8$) were enumerated with nauty and checked by a
from-scratch C program, with per-order counts matching OEIS A001174 exactly
and zero counterexamples. The CaDiCaL runs also close every shard for
$n = 15$ and $n = 16$; at the time of writing the DRAT proofs of the two
largest shards (0.6\,GB and 4.2\,GB) had not finished re-verification, so we
state $n \le 14$ as the fully machine-checked claim and report $n = 15, 16$
separately with their exact status. We found no previously published
exhaustive verification of the conjecture over all oriented graphs of a
given order, so we state the small-order record here explicitly rather than
compare against one.
\end{abstract}

\section{The conjecture and its history}\label{sec:intro}

Let $D$ be an \emph{oriented graph}: a finite directed graph with no loops,
no parallel arcs, and no digons (if $(u,v)$ is an arc then $(v,u)$ is not).
For a vertex $v$ write $\Nout(v)$ for its out-neighborhood and
$\Nsec(v)$ for its \emph{second out-neighborhood}, the set of vertices at
directed distance exactly $2$ from $v$. A vertex $v$ with
$|\Nsec(v)| \ge |\Nout(v)|$ is called a \emph{Seymour vertex}.

\begin{conjecture}[Seymour, c.~1990; see Dean and Latka~\cite{DeanLatka95}]
\label{conj:ssnc}
Every oriented graph has a Seymour vertex.
\end{conjecture}

A counterexample is an oriented graph in which \emph{every} vertex $v$
satisfies $|\Nsec(v)| < |\Nout(v)|$. The conjecture is open. The main
published partial results are:

\begin{itemize}
\item \textbf{Tournaments.} Dean's conjecture (the tournament case) was
proved by Fisher~\cite{Fisher96} using an averaging argument over
probability distributions; Havet and Thomass\'e~\cite{HavetThomasse00} gave
a combinatorial proof via median orders, producing two Seymour vertices when
the tournament has no dominated vertex.
\item \textbf{Small minimum out-degree.} Kaneko and Locke~\cite{KanekoLocke01}
proved the conjecture for oriented graphs with minimum out-degree at most
$6$. A June 2026 preprint by Sadhukhan, Sandeep, and
Sen~\cite{SadhukhanSandeepSen26} claims the case of minimum out-degree
exactly $7$, using local reductions plus computer-assisted infeasibility
checks; it has not, to our knowledge, been refereed.
\item \textbf{Approximate versions.} Chen, Shen, and Yuster~\cite{ChenShenYuster03}
proved every oriented graph has a vertex with
$|\Nsec(v)| \ge \gamma\,|\Nout(v)|$ for $\gamma = 0.657\ldots$, with later
improvements to the constant by others.
\item \textbf{Random and dense settings.} Cohn, Godbole, Wright Harkness,
and Zhang~\cite{CohnGodbole16} showed the conjecture holds asymptotically
almost surely for random tournaments and random digraphs. Espuny D\'iaz,
Gir\~ao, Granet, and Kronenberg~\cite{EspunyDiaz25} proved that a.a.s.\
every orientation of $G(n,p)$ satisfies the conjecture for fixed $p<1/2$,
and gave two reductions; their Proposition~4 shows a vertex-minimal
counterexample $D$ must have minimum out-degree greater than
$\sqrt{|V(D)|}$.
\item \textbf{Structural classes.} Fidler and Yuster~\cite{FidlerYuster07}
handled tournaments minus a star or a subtournament and digraphs of minimum
degree $|V(D)|-2$; Brantner, Brockman, Kay, and Snively~\cite{Brantner09}
handled digraphs whose underlying graph is triangle-free.
\end{itemize}

\paragraph{On the exhaustive small-order record.}
We searched for a published exhaustive verification of
Conjecture~\ref{conj:ssnc} over \emph{all} oriented graphs on $n \le k$
vertices for any $k$ and found none: neither the standard open-problem
surveys~\cite{MoharPage,WestPage} nor the papers above report one. An
earlier internal working note of ours asserted a published record of
$n \le 7$; we could not trace that assertion to any source and we withdraw
it. The correct statement appears to be that no such record was published,
presumably because checks up to $n \approx 8$ are routine with nauty. The
contribution of this note is therefore not the beating of a record but the
explicit, certificate-backed statement of one.

\begin{theorem}[main claim, machine-checked]\label{thm:main}
No oriented graph on at most $14$ vertices is a counterexample to
Conjecture~\ref{conj:ssnc}.
\end{theorem}

Combined with Kaneko and Locke~\cite{KanekoLocke01} (any counterexample has
minimum out-degree at least $7$) and the arc-counting bound of
Section~\ref{sec:shards} (minimum out-degree at most
$\lfloor (n-1)/2 \rfloor$), a hypothetical counterexample must in fact have
$n \ge 15$; Theorem~\ref{thm:main} is the unconditional statement that does
not rely on~\cite{KanekoLocke01}, and our shards re-prove the low-degree
cases for $n \le 14$ from scratch rather than citing them.

\section{Method}\label{sec:method}

The overall shape: for each order $n$, we encode ``an oriented graph on $n$
labeled vertices in which every vertex violates the Seymour inequality''
as a CNF formula, split the space into minimum out-degree shards with
symmetry breaking, obtain UNSAT verdicts with DRAT proofs from CaDiCaL
3.0.1~\cite{Cadical}, and check every proof with
drat-trim~\cite{DratTrim}. The DRAT certificates prove the CNFs
unsatisfiable; the link from the CNFs to the conjecture is the soundness
chain below, which is deliberately short and must be audited by a human.

\subsection{Encoding}\label{sec:encoding}

The encoder is \texttt{encode.py} (about 150 lines, PySAT~\cite{PySAT}
1.9). Variables: $x_{uv}$ for each ordered pair ($u \to v$ arc), and
$y_{vw}$ for each ordered pair with $v \ne w$, an \emph{upper indicator}
for $w \in \Nsec(v)$. Clauses:

\begin{enumerate}
\item $(\lnot x_{uv} \lor \lnot x_{vu})$ for all $u < v$: no digons. No
variable $x_{vv}$ exists, so loops cannot be represented.
\item $(x_{vu} \land x_{uw} \land \lnot x_{vw}) \to y_{vw}$ for all
distinct $u, v, w$: every genuine 2-path not shortcut by an arc forces
$y_{vw}$ up. Hence in every model,
$\sum_w y_{vw} \ge |\Nsec(v)|$.
\item For each $v$: a totalizer cardinality
constraint~\cite{Bailleux03} encoding
$\sum_{w} y_{vw} + \sum_{u} \lnot x_{vu} \le n-2$, which given the
previous point is equivalent to $|\Nsec(v)| \le |\Nout(v)| - 1$, i.e.,
$v$ violates the Seymour inequality.
\end{enumerate}

\emph{Completeness} (the direction that matters for a nonexistence claim):
if a counterexample $D$ on $n$ labeled vertices exists, set $x$ from the
arcs of $D$ and set $y_{vw}$ to the exact indicator of $w \in \Nsec(v)$.
Group (1) holds because $D$ is an oriented graph; group (2) holds because
any 2-path $v \to u \to w$ with no arc $v \to w$ witnesses
$w \in \Nsec(v)$; group (3) holds because every vertex of $D$ violates the
inequality. So the CNF is satisfiable, and UNSAT verdicts rule out every
labeled counterexample on $n$ vertices. Spurious $y = 1$ assignments only
tighten constraint (3), so soundness of the search is not an issue for the
nonexistence direction.

For $n \le 9$ we ran this formula over the whole space with no symmetry
breaking and a single additional constraint, minimum out-degree $\ge 1$
(the encoder's default): this excludes only graphs containing a sink, and
a sink $v$ has $|\Nsec(v)| = 0 = |\Nout(v)|$, so it is itself a Seymour
vertex and no graph containing one is a counterexample. Beyond that one
elementary observation the four instances are assumption-free: their
verified refutations alone prove Theorem~\ref{thm:main} for $n \le 9$.

\subsection{Degree shards and symmetry breaking}\label{sec:shards}

For $n \ge 10$ the whole-space formula becomes slow, so the space is split.
In a counterexample no vertex is a sink (a sink $v$ has
$|\Nsec(v)| = 0 = |\Nout(v)|$ and would be a Seymour vertex), so the
minimum out-degree $d_0$ satisfies $d_0 \ge 1$; counting arcs,
$n \, d_0 \le \sum_v \deg^+(v) = |A(D)| \le \binom{n}{2}$, so
$d_0 \le \lfloor (n-1)/2 \rfloor$. Shard $s_d$ (for
$1 \le d \le \lfloor (n-1)/2 \rfloor$) asserts:
\begin{itemize}
\item every vertex has out-degree $\ge d$ (totalizers),
\item vertex $0$ has out-degree exactly $d$ (totalizers),
\item $\Nout(0) = \{1, \dots, d\}$ (unit clauses),
\item two-block lex-leader symmetry breaking: for every adjacent
transposition $\sigma$ acting within $\{1,\dots,d\}$ or within
$\{d+1,\dots,n-1\}$, the adjacency matrix $A$ satisfies
$A \le_{\mathrm{lex}} A \circ \sigma$~\cite{Crawford96}.
\end{itemize}

\emph{Covering argument.} Suppose a counterexample $D$ on $n$ vertices
exists and let $d_0$ be its minimum out-degree, so
$1 \le d_0 \le \lfloor (n-1)/2 \rfloor$. Relabel a minimum out-degree
vertex to $0$ and its out-neighbors to $1, \dots, d_0$. The residual
relabeling freedom is exactly the subgroup
$H = \mathrm{Sym}\{1,\dots,d_0\} \times \mathrm{Sym}\{d_0+1,\dots,n-1\}$,
which preserves the pinned constraints. Choose the labeling whose adjacency
matrix is lexicographically minimal over the $H$-orbit; it satisfies
$A \le_{\mathrm{lex}} A \circ \sigma$ for \emph{every} $\sigma \in H$, in
particular for the encoded adjacent transpositions. Hence some assignment
satisfies shard $s_{d_0}$, and UNSAT verdicts on all shards of order $n$
rule out all counterexamples on $n$ vertices. Note the encoded constraints
are a \emph{subset} (adjacent transpositions only) of the full lex-leader
condition; using a subset weakens the formula and therefore cannot cause a
false UNSAT.

For $n = 10$ we additionally ran the whole space in one piece (minimum
out-degree $\ge 1$, full-lex symmetry breaking over all adjacent
transpositions of the symmetric group, which is sound because the formula
without pinning is invariant under all vertex permutations), so
Theorem~\ref{thm:main} at $n = 10$ does not depend on the shard argument
either. Cross-checks at $n = 10, 11$ with the (audited, elementary) bound
``any counterexample has minimum out-degree $\ge 3$'' agree; that lemma is
stated in the artifact repository but is not load-bearing anywhere in this
note, precisely because shards $s_1$ and $s_2$ were machine-proved
directly.

\subsection{Tournament-forced shards and Lemma T}\label{sec:lemmat}

For odd $n$, the top shard $d = (n-1)/2$ forces
$\sum_v \deg^+(v) \ge n(n-1)/2$, i.e., every pair adjacent: the graph is a
$\frac{n-1}{2}$-regular tournament. These shards are closed twice over:

\begin{enumerate}
\item By Fisher's theorem~\cite{Fisher96}: every tournament has a Seymour
vertex, so no tournament is a counterexample. This is published,
peer-reviewed mathematics and is the citation we prefer.
\item By the following elementary argument, produced in-session and
\textbf{not yet reviewed by any human}; we record it because it is short
enough to audit in one sitting.
\end{enumerate}

\begin{lemma}[Lemma T, \textbf{unreviewed}]\label{lem:t}
In a $k$-regular tournament, every vertex $v$ satisfies
$|\Nsec(v)| = |\Nout(v)| = k$, so no such tournament is a counterexample.
\end{lemma}

\begin{proof}[Proof (to be audited)]
Fix $v$ and take any $w \notin \{v\} \cup \Nout(v)$. If $w \notin
\Nsec(v)$, then no $u \in \Nout(v)$ has $u \to w$; by totality of the
tournament, $w \to u$ for all $u \in \Nout(v)$, and also $w \to v$ (since
$v \to w$ would put $w$ in $\Nout(v)$). Then $\deg^+(w) \ge k + 1$,
contradicting $k$-regularity. Hence $\Nsec(v) = V \setminus (\{v\} \cup
\Nout(v))$, which has size $n - 1 - k = k$ since $n = 2k+1$.
\end{proof}

For $n \le 14$ neither item is needed (no tournament-forced shard occurs
with an unfinished SAT run there). They matter only for the $n = 15$ and
$n = 17$ top shards discussed in Section~\ref{sec:beyond}; and for
$n = 15$ the SAT run of the tournament shard itself finished UNSAT as well,
so Lemma~T is a third line of defense there, not a dependency.

\section{Two independent pipelines}\label{sec:pipelines}

\subsection{Pipeline A: SAT with DRAT certificates}

Every instance in Table~\ref{tab:instances} was solved by CaDiCaL
3.0.1~\cite{Cadical} with binary DRAT proof output. Every completed proof
was checked with drat-trim~\cite{DratTrim} built from source (commit
\texttt{2e3b2dc}, upstream \texttt{marijnheule/drat-trim}). Of the 59 completed refutations,
57 have \texttt{s VERIFIED} receipts at freeze time; the remaining two are
the giant $n15\_s7$ and $n16\_s7$ proofs whose checks were still running
(see \texttt{receipts/SUMMARY.txt} for the final word; every receipt
generated so far reports \texttt{s VERIFIED} and none has ever failed).
The receipts were \emph{regenerated from scratch in the write-up session}
(2026-07-23) by a different agent than the one that ran the solver, and are
stored alongside this note (\texttt{receipts/}). As a negative control, a truncated copy of
one proof was rejected by drat-trim (\texttt{s NOT VERIFIED}), confirming
the checker does not rubber-stamp.

Additionally, all 63 CNF files on disk were regenerated byte-identically
from \texttt{encode.py} with their documented flags in the write-up
session, tying the artifacts that were solved to the encoder that was
audited: the four whole-space files ($n = 6,\dots,9$) with no flags, the
three cross-check files with \texttt{--mindeg} and full-lex flags, and all
56 shard files with
\texttt{--mindeg $d$ --exactdeg0 $d$ --pin-nbhd --sb-mode twoblock}.

\subsection{Pipeline B: exhaustive enumeration, $n \le 8$}

Independently of SAT, encoder, and solver: nauty's~\cite{Nauty}
\texttt{geng} generated all graphs, \texttt{directg -o} all their
orientations, and a from-scratch C checker (\texttt{check\_d6.c}, about 80
lines, reading digraph6 directly and re-checking the no-loop/no-digon
property itself) tested every oriented graph on $n \le 8$ vertices. Counts
match OEIS A001174~\cite{OEIS} (number of oriented graphs on $n$ unlabeled
vertices) exactly at every order, which is strong evidence the enumeration
is complete:

\begin{center}
\begin{tabular}{lrrrrrrrr}
\toprule
$n$ & 1 & 2 & 3 & 4 & 5 & 6 & 7 & 8 \\
\midrule
oriented graphs & 1 & 2 & 7 & 42 & 582 & 21{,}480 & 2{,}142{,}288 & 575{,}016{,}219 \\
counterexamples & 0 & 0 & 0 & 0 & 0 & 0 & 0 & 0 \\
\bottomrule
\end{tabular}
\end{center}

The $n \le 7$ runs were re-executed live during the write-up session
(seconds); the $n = 8$ run was sharded ten ways
(\texttt{geng}'s \texttt{res/mod} splitting) and the per-shard
\texttt{checked} counts sum to exactly 575{,}016{,}219. A positive control
(\texttt{check\_weak}, the same checker with the comparison deliberately
weakened) flags all 7 oriented graphs on 3 vertices, as it must.

So for $n \le 8$ the result is double-proved by disjoint toolchains
(PySAT/CaDiCaL/drat-trim versus nauty/hand-written C), and for
$n = 9,\dots,14$ it rests on Pipeline A plus the written arguments of
Section~\ref{sec:method}.

\subsection{Auxiliary lane: Cayley digraphs of small groups}

As a targeted probe beyond $n = 14$ (not part of the headline claim), all
$55$ groups of orders $12$ through $24$ were constructed and brute-force
verified to be groups, with isomorphism-type counts per order matching the
classification ($5,1,2,1,14,1,5,1,5,2,2,1,15$); for each group, every
connection set $S$ with $S \cap S^{-1} = \emptyset$ was enumerated
($1{,}505{,}160$ sets total) and the resulting Cayley digraph checked at
the identity vertex, vertex-transitivity justifying the single-vertex
check; that shortcut was itself revalidated by full all-vertex checks on
every 250th set. Zero counterexamples. This lane shares no code with either
pipeline above but has not been independently re-audited; treat it as a
search result, not a certified one.

\section{Controls}\label{sec:controls}

Beyond the negative proof-checker control and the byte-identical CNF
regeneration already described:
\begin{itemize}
\item \textbf{Slack controls (SAT expected).} Relaxing the per-vertex bound
by one (encoding $|\Nsec(v)| \le |\Nout(v)|$) must be satisfiable, since,
e.g., directed cycles and their blowups sit exactly on the boundary. This
was confirmed at $n = 3$ (plain) and, critically, at $n = 6$ and $n = 12$
\emph{through the full shard machinery} (pinning, exact degree, two-block
lex): the shard path is not accidentally over-constrained. The $n = 12$
control model was decoded back into an explicit digraph and checked to be a
genuine boundary witness, validating the totalizer calibration end to end.
\item \textbf{Shard UNSAT control.} A shard instance at $n = 7$, where
Pipeline B independently proves no counterexample exists, is UNSAT as
expected. All four control instances were re-solved from the stored CNF
files during the write-up session with the expected verdicts
(SAT, SAT, UNSAT, SAT).
\item \textbf{Solver agreement.} Spot-checked shards were re-solved with
Glucose 4.2 and MiniSat 2.2 with agreeing verdicts (no proofs logged from
those; CaDiCaL's DRAT proofs are the receipts).
\end{itemize}

\section{Results}\label{sec:results}

Table~\ref{tab:instances} lists every SAT instance with its DIMACS size,
CaDiCaL wall time, DRAT proof size, verdict, and this session's drat-trim
receipt. The table is generated mechanically from the artifacts by
\texttt{gen\_tables.py}; no number in it was typed by hand.
Wall times are from a heavily shared machine and should be read as upper
bounds. Shard counts per order equal $\lfloor (n-1)/2 \rfloor$ exactly:
$4$ shards at $n=10$, then $5,5,6,6,7,7,8,8$ through $n=18$.

% auto-generated by gen_tables.py; do not edit by hand
{\footnotesize
\setlength{\tabcolsep}{4pt}
\begin{longtable}{l l r r r r l l}
\caption{All SAT instances. Columns: DIMACS size, CaDiCaL wall time (seconds), DRAT proof size, solver verdict, and the drat-trim receipt regenerated in the write-up session (2026-07-23).}\label{tab:instances}\\
\toprule
instance & meaning & vars & clauses & time & proof & verdict & drat-trim \\
\midrule\endfirsthead
\toprule instance & meaning & vars & clauses & time & proof & verdict & drat-trim \\ \midrule\endhead
\bottomrule\endfoot
\texttt{n6} & $n{=}6$ whole space, $\delta^+{\ge}1$, no SB & 264 & 621 & $<$0.1 & 14\,KB & UNSAT & \textsc{verified} \\
\texttt{n7} & $n{=}7$ whole space, $\delta^+{\ge}1$, no SB & 392 & 1015 & 0.4 & 137\,KB & UNSAT & \textsc{verified} \\
\texttt{n8} & $n{=}8$ whole space, $\delta^+{\ge}1$, no SB & 544 & 1540 & 10 & 1.9\,MB & UNSAT & \textsc{verified} \\
\texttt{n9} & $n{=}9$ whole space, $\delta^+{\ge}1$, no SB & 720 & 2214 & 279 & 31.2\,MB & UNSAT & \textsc{verified} \\
\texttt{n10\_d1\_sb} & $n{=}10$ whole space, $\delta^+{\ge}1$, full lex & 1246 & 4902 & 27 & 5.7\,MB & UNSAT & \textsc{verified} \\
\texttt{n10\_d3\_sb} & $n{=}10$, $\delta^+{\ge}3$ (lemma), full lex & 1536 & 5552 & 15 & 3.6\,MB & UNSAT & \textsc{verified} \\
\texttt{n11\_d3\_sb} & $n{=}11$, $\delta^+{\ge}3$ (lemma), full lex & 1942 & 7264 & 503 & 55.5\,MB & UNSAT & \textsc{verified} \\
\texttt{n10\_s1} & $n{=}10$, shard $d_0{=}1$ & 1207 & 4571 & $<$0.1 & 11\,KB & UNSAT & \textsc{verified} \\
\texttt{n10\_s2} & $n{=}10$, shard $d_0{=}2$ & 1497 & 5221 & $<$0.1 & 2\,KB & UNSAT & \textsc{verified} \\
\texttt{n10\_s3} & $n{=}10$, shard $d_0{=}3$ & 1497 & 5221 & $<$0.1 & 4\,KB & UNSAT & \textsc{verified} \\
\texttt{n10\_s4} & $n{=}10$, shard $d_0{=}4$ & 1497 & 5221 & $<$0.1 & 4\,KB & UNSAT & \textsc{verified} \\
\texttt{n11\_s1} & $n{=}11$, shard $d_0{=}1$ & 1526 & 6031 & $<$0.1 & 14\,KB & UNSAT & \textsc{verified} \\
\texttt{n11\_s2} & $n{=}11$, shard $d_0{=}2$ & 1900 & 6900 & $<$0.1 & 2\,KB & UNSAT & \textsc{verified} \\
\texttt{n11\_s3} & $n{=}11$, shard $d_0{=}3$ & 1900 & 6900 & $<$0.1 & 5\,KB & UNSAT & \textsc{verified} \\
\texttt{n11\_s4} & $n{=}11$, shard $d_0{=}4$ & 1900 & 6900 & $<$0.1 & 7\,KB & UNSAT & \textsc{verified} \\
\texttt{n11\_s5} & $n{=}11$, shard $d_0{=}5$ & 1900 & 6900 & $<$0.1 & 30\,KB & UNSAT & \textsc{verified} \\
\texttt{n12\_s1} & $n{=}12$, shard $d_0{=}1$ & 1881 & 7747 & $<$0.1 & 17\,KB & UNSAT & \textsc{verified} \\
\texttt{n12\_s2} & $n{=}12$, shard $d_0{=}2$ & 2349 & 8875 & $<$0.1 & 2\,KB & UNSAT & \textsc{verified} \\
\texttt{n12\_s3} & $n{=}12$, shard $d_0{=}3$ & 2349 & 8875 & $<$0.1 & 5\,KB & UNSAT & \textsc{verified} \\
\texttt{n12\_s4} & $n{=}12$, shard $d_0{=}4$ & 2349 & 8875 & $<$0.1 & 43\,KB & UNSAT & \textsc{verified} \\
\texttt{n12\_s5} & $n{=}12$, shard $d_0{=}5$ & 2349 & 8875 & 0.2 & 111\,KB & UNSAT & \textsc{verified} \\
\texttt{n13\_s1} & $n{=}13$, shard $d_0{=}1$ & 2272 & 9737 & $<$0.1 & 22\,KB & UNSAT & \textsc{verified} \\
\texttt{n13\_s2} & $n{=}13$, shard $d_0{=}2$ & 2844 & 11167 & $<$0.1 & 3\,KB & UNSAT & \textsc{verified} \\
\texttt{n13\_s3} & $n{=}13$, shard $d_0{=}3$ & 2844 & 11167 & $<$0.1 & 6\,KB & UNSAT & \textsc{verified} \\
\texttt{n13\_s4} & $n{=}13$, shard $d_0{=}4$ & 2844 & 11167 & $<$0.1 & 8\,KB & UNSAT & \textsc{verified} \\
\texttt{n13\_s5} & $n{=}13$, shard $d_0{=}5$ & 2844 & 11167 & 0.7 & 149\,KB & UNSAT & \textsc{verified} \\
\texttt{n13\_s6} & $n{=}13$, shard $d_0{=}6$ & 2844 & 11167 & 10.0 & 1.5\,MB & UNSAT & \textsc{verified} \\
\texttt{n14\_s1} & $n{=}14$, shard $d_0{=}1$ & 2699 & 12019 & $<$0.1 & 26\,KB & UNSAT & \textsc{verified} \\
\texttt{n14\_s2} & $n{=}14$, shard $d_0{=}2$ & 3385 & 13797 & $<$0.1 & 3\,KB & UNSAT & \textsc{verified} \\
\texttt{n14\_s3} & $n{=}14$, shard $d_0{=}3$ & 3385 & 13797 & $<$0.1 & 5\,KB & UNSAT & \textsc{verified} \\
\texttt{n14\_s4} & $n{=}14$, shard $d_0{=}4$ & 3385 & 13797 & $<$0.1 & 41\,KB & UNSAT & \textsc{verified} \\
\texttt{n14\_s5} & $n{=}14$, shard $d_0{=}5$ & 3385 & 13797 & 2.3 & 486\,KB & UNSAT & \textsc{verified} \\
\texttt{n14\_s6} & $n{=}14$, shard $d_0{=}6$ & 3385 & 13797 & 41 & 6.4\,MB & UNSAT & \textsc{verified} \\
\texttt{n15\_s1} & $n{=}15$, shard $d_0{=}1$ & 3162 & 14611 & $<$0.1 & 31\,KB & UNSAT & \textsc{verified} \\
\texttt{n15\_s2} & $n{=}15$, shard $d_0{=}2$ & 3972 & 16786 & $<$0.1 & 3\,KB & UNSAT & \textsc{verified} \\
\texttt{n15\_s3} & $n{=}15$, shard $d_0{=}3$ & 3972 & 16786 & $<$0.1 & 6\,KB & UNSAT & \textsc{verified} \\
\texttt{n15\_s4} & $n{=}15$, shard $d_0{=}4$ & 3972 & 16786 & $<$0.1 & 73\,KB & UNSAT & \textsc{verified} \\
\texttt{n15\_s5} & $n{=}15$, shard $d_0{=}5$ & 3972 & 16786 & 2.8 & 568\,KB & UNSAT & \textsc{verified} \\
\texttt{n15\_s6} & $n{=}15$, shard $d_0{=}6$ & 3972 & 16786 & 51 & 9.9\,MB & UNSAT & \textsc{verified} \\
\texttt{n15\_s7} & $n{=}15$, shard $d_0{=}7$ & 3972 & 16786 & 4643 & 568.4\,MB & UNSAT & pending \\
\texttt{n16\_s1} & $n{=}16$, shard $d_0{=}1$ & 3661 & 17531 & $<$0.1 & 36\,KB & UNSAT & \textsc{verified} \\
\texttt{n16\_s2} & $n{=}16$, shard $d_0{=}2$ & 4605 & 20155 & $<$0.1 & 3\,KB & UNSAT & \textsc{verified} \\
\texttt{n16\_s3} & $n{=}16$, shard $d_0{=}3$ & 4605 & 20155 & $<$0.1 & 7\,KB & UNSAT & \textsc{verified} \\
\texttt{n16\_s4} & $n{=}16$, shard $d_0{=}4$ & 4605 & 20155 & $<$0.1 & 85\,KB & UNSAT & \textsc{verified} \\
\texttt{n16\_s5} & $n{=}16$, shard $d_0{=}5$ & 4605 & 20155 & 3.9 & 755\,KB & UNSAT & \textsc{verified} \\
\texttt{n16\_s6} & $n{=}16$, shard $d_0{=}6$ & 4605 & 20155 & 116 & 24.2\,MB & UNSAT & \textsc{verified} \\
\texttt{n16\_s7} & $n{=}16$, shard $d_0{=}7$ & 4605 & 20155 & 19080 & 3.96\,GB & UNSAT & pending \\
\texttt{n17\_s1} & $n{=}17$, shard $d_0{=}1$ & 4196 & 20797 & $<$0.1 & 42\,KB & UNSAT & \textsc{verified} \\
\texttt{n17\_s2} & $n{=}17$, shard $d_0{=}2$ & 5284 & 23925 & $<$0.1 & 4\,KB & UNSAT & \textsc{verified} \\
\texttt{n17\_s3} & $n{=}17$, shard $d_0{=}3$ & 5284 & 23925 & $<$0.1 & 8\,KB & UNSAT & \textsc{verified} \\
\texttt{n17\_s4} & $n{=}17$, shard $d_0{=}4$ & 5284 & 23925 & $<$0.1 & 49\,KB & UNSAT & \textsc{verified} \\
\texttt{n17\_s5} & $n{=}17$, shard $d_0{=}5$ & 5284 & 23925 & 2.6 & 774\,KB & UNSAT & \textsc{verified} \\
\texttt{n17\_s6} & $n{=}17$, shard $d_0{=}6$ & 5284 & 23925 & 84 & 18.6\,MB & UNSAT & \textsc{verified} \\
\texttt{n17\_s7} & $n{=}17$, shard $d_0{=}7$ & 5284 & 23925 & -- & -- & running & -- \\
\texttt{n17\_s8} & $n{=}17$, shard $d_0{=}8$ & 5284 & 23925 & -- & -- & running & -- \\
\texttt{n18\_s1} & $n{=}18$, shard $d_0{=}1$ & 4804 & 24464 & $<$0.1 & 48\,KB & UNSAT & \textsc{verified} \\
\texttt{n18\_s2} & $n{=}18$, shard $d_0{=}2$ & 6064 & 28172 & $<$0.1 & 4\,KB & UNSAT & \textsc{verified} \\
\texttt{n18\_s3} & $n{=}18$, shard $d_0{=}3$ & 6064 & 28172 & $<$0.1 & 8\,KB & UNSAT & \textsc{verified} \\
\texttt{n18\_s4} & $n{=}18$, shard $d_0{=}4$ & 6064 & 28172 & $<$0.1 & 110\,KB & UNSAT & \textsc{verified} \\
\texttt{n18\_s5} & $n{=}18$, shard $d_0{=}5$ & 6064 & 28172 & 2.9 & 927\,KB & UNSAT & \textsc{verified} \\
\texttt{n18\_s6} & $n{=}18$, shard $d_0{=}6$ & 6064 & 28172 & 113 & 25.2\,MB & UNSAT & \textsc{verified} \\
\texttt{n18\_s7} & $n{=}18$, shard $d_0{=}7$ & 6064 & 28172 & -- & -- & running & -- \\
\texttt{n18\_s8} & $n{=}18$, shard $d_0{=}8$ & 6064 & 28172 & -- & -- & running & -- \\
\end{longtable}
}


\subsection{Fully verified claim}

For $n \le 14$: every listed instance is UNSAT with a drat-trim-verified
DRAT proof. By the completeness argument of Section~\ref{sec:encoding}
(for $n \le 9$; the $n = 10$ whole-space instance additionally uses the
full-lex soundness argument of Section~\ref{sec:shards}) and the covering
argument of Section~\ref{sec:shards} (for the sharded orders), this proves
Theorem~\ref{thm:main}. For $n \le 8$, Pipeline B independently confirms
it.

\subsection{Beyond 14: exact status at write-up time}\label{sec:beyond}

\begin{itemize}
\item $\mathbf{n = 15}$: all $7$ shards UNSAT by CaDiCaL (top shard
$s_7$: $4{,}643$\,s wall, $0.60$\,GB proof). Shards $s_1$ through $s_6$
have verified DRAT proofs; verification of the $s_7$ proof was still
running when this note was frozen (its receipt, if present in
\texttt{receipts/}, supersedes this sentence). Independently of that one
certificate, shard $s_7$ forces a $7$-regular tournament and is closed by
Fisher's theorem~\cite{Fisher96}, and also by the unreviewed
Lemma~\ref{lem:t}. We therefore regard $n = 15$ as established modulo one
of: (a) the pending big-proof verification, or (b) reliance on
Fisher~\cite{Fisher96}, and we deliberately keep it out of the headline.
\item $\mathbf{n = 16}$: all $7$ shards UNSAT by CaDiCaL ($s_7$:
$19{,}080$\,s wall, $4.23$\,GB proof); $s_1$ through $s_6$ verified;
verification of the $s_7$ proof still running. No tournament argument
applies at even $n$, so $n = 16$ waits on that verification.
\item $\mathbf{n = 17, 18}$: shards $s_1$ through $s_6$ UNSAT and
verified. The near-tournament shards ($n17\_s7$, $n17\_s8$, $n18\_s7$,
$n18\_s8$) were still being solved at write-up time ($>5$ hours each).
$n17\_s8$ forces an $8$-regular tournament and is closed by
Fisher~\cite{Fisher96} (or Lemma~\ref{lem:t}) regardless of its SAT run;
$n17\_s7$, $n18\_s7$, $n18\_s8$ are genuinely open runs. So $n = 17, 18$
are \emph{not} claimed.
\end{itemize}

\section{What an external reviewer must audit}\label{sec:audit}

This note was produced by an automated pipeline. Before any submission, a
human must independently check, at minimum:

\begin{enumerate}
\item \textbf{The soundness chain (Section~\ref{sec:method}).} The DRAT
certificates prove only that specific CNF files are unsatisfiable. The
encoding-completeness argument, the sink and arc-count bounds on $d_0$,
and the two-block lex-leader covering argument are pen-and-paper claims
living outside all certificates. They are short; audit every line.
\item \textbf{\texttt{encode.py} against Section~\ref{sec:encoding}.}
About 150 lines. Confirm the clause groups match the stated encoding and
that PySAT's totalizer (\texttt{CardEnc.atmost}, \texttt{EncType.totalizer})
is used correctly (auxiliary variables must not be shared across
constraints incorrectly; the slack controls test this but do not prove it).
\item \textbf{Lemma~\ref{lem:t}}, or simply replace every use of it by the
citation to Fisher~\cite{Fisher96} and delete it.
\item \textbf{Re-verification of all certificates on independent hardware},
rebuilding CaDiCaL and drat-trim from upstream sources. The full artifact
set is about 28\,GB; SHA-256 checksums of every CNF, every completed proof,
and the three load-bearing source files are in
\texttt{receipts/SHA256SUMS.txt}. Ideally re-check with a second,
independently implemented DRAT checker (e.g., the formally verified
\texttt{cake\_lpr}).
\item \textbf{Pipeline B.} Read \texttt{check\_d6.c} (about 80 lines,
including its own digraph6 parser and digon check); confirm the
\texttt{geng}/\texttt{directg -o} invocations enumerate exactly the
oriented graphs; re-derive the OEIS comparison.
\item \textbf{Provenance discipline.} Only \texttt{encode.py},
\texttt{check\_d6.c}, the CNF/DRAT artifacts, and the receipts are
load-bearing. Earlier search code from the same overnight runs (local
search, circulant search, the round-1 SAT stack) is untrusted by policy
and cited nowhere in the argument. Confirm nothing in the chain silently
depends on it.
\item \textbf{Literature claims.} Re-check the history of
Section~\ref{sec:intro} against the primary sources, in particular the
exact statement of Kaneko--Locke and the status of the 2026 minimum
out-degree 7 preprint, and re-run the search for any published exhaustive
small-order verification we may have missed.
\item \textbf{The $n = 15, 16$ upgrades.} If the two large proof
verifications complete with \texttt{s VERIFIED} receipts, the headline can
move to $n \le 16$; that decision, and the wording, should be a human's.
\end{enumerate}

\section*{Reproduction appendix}

All commands assume macOS with Homebrew; everything is CPU-only. On a
shared machine, prefix heavy runs with \texttt{nice -n 15}.

\begin{enumerate}
\item Toolchain:
\begin{verbatim}
brew install cadical nauty python
pip install python-sat
git clone https://github.com/marijnheule/drat-trim && make -C drat-trim
\end{verbatim}
\item Whole-space instances (assumption-free), $n = 6, \dots, 9$:
\begin{verbatim}
python3 encode.py N --out nN.cnf
cadical -q nN.cnf nN.drat          # expect: s UNSATISFIABLE
drat-trim/drat-trim nN.cnf nN.drat # expect: s VERIFIED
\end{verbatim}
\item Shard instances, $n = 10, \dots, 18$, $d = 1, \dots,
\lfloor (n-1)/2 \rfloor$:
\begin{verbatim}
python3 encode.py N --mindeg D --exactdeg0 D --pin-nbhd \
    --sb-mode twoblock --out nN_sD.cnf
cadical -q nN_sD.cnf nN_sD.drat
drat-trim/drat-trim nN_sD.cnf nN_sD.drat
\end{verbatim}
\item Whole-space cross-checks at $n = 10, 11$:
\texttt{--mindeg 1 --sb-mode full} (and \texttt{--mindeg 3} for the
lemma-conditioned variants).
\item Controls: \texttt{--slack} variants must come back SAT; truncating
any \texttt{.drat} must make drat-trim report \texttt{s NOT VERIFIED}.
\item Pipeline B ($n \le 8$; shard $n = 8$ via \texttt{res/mod}):
\begin{verbatim}
cc -O2 -o check_d6 check_d6.c
geng -q N | directg -o -q | ./check_d6
geng -q 8 r/10 | directg -o -q | ./check_d6  # r = 0..9, sum counts
\end{verbatim}
\item Regenerate Table~\ref{tab:instances} from the artifacts, and
re-verify every stored proof:
\begin{verbatim}
python3 gen_tables.py
sh verify_all_drat.sh
\end{verbatim}
\end{enumerate}

\begin{thebibliography}{99}

\bibitem{DeanLatka95} N.~Dean and B.~J.~Latka, Squaring the tournament: an
open problem, \emph{Congressus Numerantium} \textbf{109} (1995), 73--80.

\bibitem{Fisher96} D.~C.~Fisher, Squaring a tournament: a proof of Dean's
conjecture, \emph{Journal of Graph Theory} \textbf{23} (1996), 43--48.

\bibitem{HavetThomasse00} F.~Havet and S.~Thomass\'e, Median orders of
tournaments: a tool for the second neighborhood problem and Sumner's
conjecture, \emph{Journal of Graph Theory} \textbf{35} (2000), 244--256.

\bibitem{KanekoLocke01} Y.~Kaneko and S.~C.~Locke, The minimum degree
approach for Paul Seymour's distance 2 conjecture, \emph{Congressus
Numerantium} \textbf{148} (2001), 201--206.

\bibitem{ChenShenYuster03} G.~Chen, J.~Shen, and R.~Yuster, Second
neighborhood via first neighborhood in digraphs, \emph{Annals of
Combinatorics} \textbf{7} (2003), 15--20.

\bibitem{FidlerYuster07} D.~Fidler and R.~Yuster, Remarks on the second
neighborhood problem, \emph{Journal of Graph Theory} \textbf{55} (2007),
208--220.

\bibitem{Brantner09} J.~Brantner, G.~Brockman, B.~Kay, and E.~Snively,
Contributions to Seymour's second neighborhood conjecture, \emph{Involve}
\textbf{2} (2009), 387--395.

\bibitem{CohnGodbole16} Z.~Cohn, A.~Godbole, E.~Wright Harkness, and
Y.~Zhang, The number of Seymour vertices in random tournaments and
digraphs, \emph{Graphs and Combinatorics} \textbf{32} (2016), 1805--1816.

\bibitem{EspunyDiaz25} A.~Espuny D\'iaz, A.~Gir\~ao, B.~Granet, and
G.~Kronenberg, Seymour's second neighbourhood conjecture: random graphs
and reductions, \emph{Random Structures \& Algorithms} (2025),
arXiv:2403.02842.

\bibitem{SadhukhanSandeepSen26} A.~Sadhukhan, R.~B.~Sandeep, and S.~Sen,
A proof of Seymour's second neighborhood conjecture for oriented graphs
with minimum out-degree equal to 7, arXiv:2606.30588 (2026). Preprint,
unrefereed as far as we know.

\bibitem{MoharPage} B.~Mohar, Seymour's Second Neighborhood Conjecture,
problem page,
\url{https://www.sfu.ca/~mohar/Problems/P0601_SeymourSecondNbdConj.html},
accessed 2026-07-23.

\bibitem{WestPage} D.~B.~West, Seymour's 2nd Neighborhood Conjecture, open
problem page, \url{https://dwest.web.illinois.edu/openp/2ndnbhd.html},
accessed 2026-07-23.

\bibitem{Cadical} A.~Biere, K.~Fazekas, M.~Fleury, and M.~Heisinger,
CaDiCaL, Kissat, Paracooba, Plingeling and Treengeling entering the SAT
Competition 2020, in \emph{Proc.\ SAT Competition 2020}, University of
Helsinki, 2020. Version used here: CaDiCaL 3.0.1.

\bibitem{DratTrim} N.~Wetzler, M.~J.~H.~Heule, and W.~A.~Hunt~Jr.,
DRAT-trim: efficient checking and trimming using expressive clausal
proofs, in \emph{Theory and Applications of Satisfiability Testing (SAT
2014)}, LNCS 8561, Springer, 2014, 422--429.

\bibitem{PySAT} A.~Ignatiev, A.~Morgado, and J.~Marques-Silva, PySAT: a
Python toolkit for prototyping with SAT oracles, in \emph{SAT 2018}, LNCS
10929, Springer, 2018, 428--437.

\bibitem{Bailleux03} O.~Bailleux and Y.~Boufkhad, Efficient CNF encoding
of Boolean cardinality constraints, in \emph{Principles and Practice of
Constraint Programming (CP 2003)}, LNCS 2833, Springer, 2003, 108--122.

\bibitem{Crawford96} J.~Crawford, M.~Ginsberg, E.~Luks, and A.~Roy,
Symmetry-breaking predicates for search problems, in \emph{Proc.\ 5th
International Conference on Principles of Knowledge Representation and
Reasoning (KR'96)}, 1996, 148--159.

\bibitem{Nauty} B.~D.~McKay and A.~Piperno, Practical graph isomorphism,
II, \emph{Journal of Symbolic Computation} \textbf{60} (2014), 94--112.

\bibitem{OEIS} OEIS Foundation Inc., entry A001174 (number of oriented
graphs on $n$ nodes), \emph{The On-Line Encyclopedia of Integer
Sequences}, \url{https://oeis.org/A001174}.

\end{thebibliography}

\end{document}
